%% ============================================================
\section{Background and Related Work}
\label{sec:background}
%% ============================================================

\subsection{Collaborative Development Environments}

The concept of a \emph{Collaborative Development Environment} (CDE) predates
the web era; early work by Dewan and Choudhary~\cite{dewan1995flexible}
described multi-user programming environments based on shared data
structures. The advent of cloud computing shifted CDEs to browser-based
deployments, yielding a diverse ecosystem of tools.

\textbf{Commercial platforms.}
GitHub Codespaces~\cite{githubcodespaces2023} provides cloud-hosted VS Code
instances accessible via any browser, with collaborative features mediated by
the Live Share protocol. CodeSandbox~\cite{codesandbox2023} offers a
browser-based sandbox with near-instant forking and preview deployment. Replit
Teams~\cite{replit2024} supports real-time collaborative editing with built-in
execution and an in-IDE AI assistant. StackBlitz~\cite{stackblitz2023}
leverages WebContainers---running Node.js natively in the browser---to
eliminate server-side execution dependencies entirely. Despite their maturity,
these platforms are either proprietary (precluding scholarly replication),
lack granular RBAC, or decouple AI assistance from the shared workspace.

\textbf{Research prototypes.}
Copdit~\cite{cheng2017codeit} explored peer-programming support in IDEs,
identifying communication friction as the dominant bottleneck. CodeTogether
(formerly eXprogramming)~\cite{codetogether2021} demonstrated that
pair-programming productivity improves significantly when the \emph{driver}
and \emph{navigator} share a live cursor and selection state. Sousa
\etal~\cite{sousa2022realtime} proposed a CRDT-based collaborative editor
achieving conflict-free merge in $O(1)$ time per operation, outperforming
OT-based approaches for high-concurrency scenarios.

\subsection{Real-Time Synchronisation Protocols}

Real-time collaborative text editing is a well-studied problem with two
dominant algorithmic families:

\textbf{Operational Transformation (OT).}
Introduced by Ellis and Gibbs~\cite{ellis1989concurrency} and subsequently
formalised by Sun \etal~\cite{sun1998operational}, OT transforms concurrent
operations to preserve document consistency. Google Docs employs a
server-mediated OT variant~\cite{googledocs2010}. OT requires a central
authority to order concurrent edits, which simplifies implementation but
introduces a single point of contention.

\textbf{Conflict-Free Replicated Data Types (CRDTs).}
Proposed by Shapiro \etal~\cite{shapiro2011conflict}, CRDTs guarantee
\emph{strong eventual consistency} without central coordination. Logoot
\cite{weiss2009logoot} and LSEQ~\cite{nedelec2013lseq} adapt CRDT principles
specifically to ordered sequence types (text documents), enabling fully
peer-to-peer collaborative editing. Automerge~\cite{kleppmann2019local} and
Yjs~\cite{nicolaescu2015yjs} provide production-grade JavaScript CRDT
implementations used by tools including Notion and Liveblocks.

\echoray{} adopts a \emph{server-relay} broadcast model via Socket.io rather
than full OT or CRDT, serialising file-tree updates through the server to
ensure authoritative consistency. This design trades theoretical generality
for implementation simplicity and straightforward auditability, a trade-off we
examine empirically in \cref{sec:evaluation}.

\subsection{WebSocket and Event-Driven Architectures}

The WebSocket protocol (RFC~6455~\cite{rfc6455}) provides full-duplex
communication over a single TCP connection with minimal framing overhead.
Socket.io~\cite{socketio2024} builds upon WebSocket (falling back to HTTP
long-polling when necessary) with a high-level event API, automatic
reconnection, and \emph{room}-based multicast, making it the de-facto choice
for real-time web applications. Comparative benchmarks by
Ramirez~\etal~\cite{ramirez2023benchmark} showed Socket.io achieving
sub-50\,ms round-trip latency under 100 concurrent connections on commodity
hardware, a result consistent with our own measurements in \cref{sec:evaluation}.

\subsection{AI-Assisted Programming}

The application of machine learning to source code has a rich history spanning
statistical language models~\cite{hindle2012naturalness}, sequence-to-sequence
architectures~\cite{bahdanau2015neural}, and, most recently, the Transformer
model~\cite{vaswani2017attention}.

\textbf{Large language models for code.}
Codex~\cite{chen2021evaluating}, the model underpinning GitHub Copilot,
demonstrated that fine-tuning GPT-3 on public code repositories yields
remarkable code-completion capabilities. Subsequent models---StarCoder
\cite{li2023starcoder}, Code Llama~\cite{roziere2023codellama}, and Google's
Gemini~\cite{team2023gemini}---extended these capabilities to multi-file
contexts, executable code generation, and structured output. Gemini~2.5~Flash,
the LLM powering \echoray{}'s AI assistant, is specifically optimised for
low-latency inference while maintaining strong reasoning capability, making it
well-suited to interactive programming assistance.

\textbf{Empirical studies of AI assistance.}
Peng \etal~\cite{peng2023impact} found that developers using AI coding
assistants completed a web-server task~55.8\% faster. Mozannar
\etal~\cite{mozannar2024reading} studied how developers read and act on
AI-generated code, finding that developers accept suggestions with low
scrutiny---motivating structured, JSON-formatted AI output (as in \echoray{})
to make AI contributions auditable. Vaithilingam \etal~\cite{vaithilingam2022expectation}
identified \emph{context sharing} and \emph{prompt visibility} as barriers to
effective AI assistance in team settings, exactly the gap our
\emph{AI-broadcast} pattern addresses.

\subsection{Authentication and Session Management}

JSON Web Tokens (JWT)~\cite{jones2015rfc7519} are widely used for stateless
authentication in REST APIs. However, their stateless nature complicates token
revocation: once issued, a JWT remains valid until expiry unless an external
revocation list is maintained. OWASP~\cite{owasp2024jwt} recommends
server-side token blacklisting backed by a cache such as Redis for scenarios
requiring immediate revocation (e.g., logout, password change, session
hijacking response). \echoray{} implements this pattern, storing blacklisted
tokens with a TTL equal to the token's remaining validity period.

\subsection{Role-Based Access Control in Collaborative Systems}

Ferraiolo and Kuhn~\cite{ferraiolo1992role} formalised RBAC as a policy model
in which permissions are assigned to roles and roles are assigned to users.
For collaborative software platforms, granular RBAC is essential: a project
owner (``leader'') must be able to add or remove collaborators and delete the
project, while ordinary members should have read/write access to project
artefacts without administrative privileges. Sandhu \etal's RBAC96
model~\cite{sandhu1996role} extended the original framework with role
hierarchies and constraints, providing a formal foundation for the simpler
two-role model implemented in \echoray{}.

\subsection{Summary and Positioning of \echoray{}}

\Cref{tab:related} summarises the most relevant tools along the five
dimensions that \echoray{} targets. Unlike existing systems, \echoray{}
simultaneously achieves open-source accessibility, AI-broadcast collaboration,
fine-grained RBAC, browser-native execution, and comprehensive security
mechanisms in a single, cohesive platform.

\begin{table}[H]
  \centering
  \caption{Comparison of collaborative development platforms across key
    dimensions. \checkmark = fully supported, $\sim$ = partially supported,
    $\times$ = not supported.}
  \label{tab:related}
  \setlength{\tabcolsep}{4pt}
  \begin{tabular}{lccccc}
    \toprule
    \textbf{Platform} & \textbf{Open} & \textbf{AI} & \textbf{AI} & \textbf{RBAC} & \textbf{In-browser} \\
                      & \textbf{Source} & \textbf{Assist} & \textbf{Broadcast} & & \textbf{Execution} \\
    \midrule
    GitHub Codespaces  & $\times$ & \checkmark & $\times$ & $\sim$ & $\times$ \\
    Replit Teams       & $\times$ & \checkmark & $\times$ & $\sim$ & \checkmark \\
    CodeSandbox        & $\sim$   & \checkmark & $\times$ & $\times$ & \checkmark \\
    StackBlitz         & $\sim$   & $\sim$     & $\times$ & $\times$ & \checkmark \\
    VS Live Share      & $\times$ & \checkmark & $\times$ & $\times$ & $\times$ \\
    Sousa \etal~(2022) & \checkmark & $\times$ & $\times$ & $\times$ & $\times$ \\
    \textbf{\echoray{}} & \checkmark & \checkmark & \checkmark & \checkmark & \checkmark \\
    \bottomrule
  \end{tabular}
\end{table}
